\chapter*{Anexos del subdocumento 6}
\addcontentsline{toc}{chapter}{Anexos del Subdocumento 6}

Estos anexos presentan los registros de interesados y adquisiciones asociados al proyecto, que complementan la información resumida en el Subdocumento 6.

\section*{Anexo 6.A \textemdash{} Interesados del proyecto}
\addcontentsline{toc}{section}{Anexo 6.A \textemdash{} Interesados del proyecto}

\setcounter{table}{0}\renewcommand{\thetable}{6.A.\arabic{table}}


El registro de interesados del proyecto (paquete 1.4.1 del Formulario T-14) usa los mismos 19 actores, con los mismos niveles de influencia e interés, que la matriz del Subdocumento 2 (Tabla 2.2 y Anexo 2.3) y la estrategia de participación del Subdocumento 3 (Anexo 3.I, Tabla 3.A.12). Este anexo agrega lo que la metodología de gestión necesita para operar ese registro: la estrategia que corresponde a cada actor según su celda de la matriz, el momento del proyecto en que participa, el responsable de LafroX y el canal por el que se le informa y consulta.

La influencia indica cuánto puede cambiar un actor el resultado del proyecto, y el interés, cuánto cambia el proyecto su trabajo. Ningún actor tiene interés bajo, porque el proyecto toca a los 19. De la matriz se derivan cuatro estrategias:
\begin{itemize}
  \item \textbf{Decidir con la dirección} (influencia muy alta e interés muy alto): participa en el Comité Ejecutivo y aprueba prioridades, cambios y cierres.
  \item \textbf{Co-diseñar las reglas} (influencia alta e interés muy alto): aporta las reglas que hoy no están escritas y valida los cambios antes de aplicarlos.
  \item \textbf{Consultar y capacitar en terreno} (influencia alta e interés alto): prueba los cambios en su propio puesto o ruta y se certifica antes de cada ola.
  \item \textbf{Informar y acordar} (influencia muy alta o media e interés alto, sin participación en la operación diaria): recibe información oportuna y firma los acuerdos formales que el proyecto necesita, sin exigirle tecnología nueva cuando no corresponde.
\end{itemize}

\begin{tablalafrox}{Registro de interesados del proyecto. Fuente: elaboración propia a partir de (Subdocumento 2, Tabla 2.2 y Anexo 2.3) y (Subdocumento 3, Anexo 3.I, Tabla 3.A.12).}{tab:tabla-interesados}%
  {F{0.20} F{0.13} Y F{0.15} F{0.17}}%
  {\cab{Actor} & \cab{Influencia / interés} & \cab{Estrategia, participación y momento} & \cab{Responsable LafroX} & \cab{Canal principal}}%
  [Registro de interesados del proyecto]
  Gerenta General & Muy alta / muy alto & Decidir: preside el Comité Ejecutivo mensual, arbitra el reparto entre etapas (mes 1) y aprueba el cierre de cada marcha blanca. & Jefe de Proyecto & Comité Ejecutivo; informe mensual \\
  Gerente Comercial & Muy alta / muy alto & Decidir: valida la promesa de entrega y las reglas comerciales (mes 2) y los cambios que afecten la venta. & Líder de Implantación & Comité Ejecutivo; talleres de levantamiento \\
  Gerente de Finanzas & Muy alta / muy alto & Decidir: valida la rendición por camión (mes 1) y la metodología de costo de servir antes del mes 13. & Líder de Datos & Comité Ejecutivo; talleres de levantamiento \\
  Gerente de Operaciones & Alta / muy alto & Co-diseñar: valida olas, reversión y dotación de estabilización antes de cada ola. & Jefe de Proyecto & Comité de Proyecto; acta de avance de ola \\
  Jefa de Calidad & Alta / muy alto & Co-diseñar: define la regla térmica por producto (mes 2), decide liberar, bloquear o rechazar y conduce el simulacro de retiro. & Líder de Datos & Talleres de reglas; Comité de Proyecto \\
  Jefa de Bodega & Alta / muy alto & Co-diseñar: valida los flujos de bodega en el turno de noche y acompaña su marcha blanca. & Líder de Implantación & Sesión de turno nocturno; Comité de Proyecto \\
  Jefe de TI & Alta / muy alto & Co-diseñar: valida las interfaces (meses 1 a 4) y recibe la transferencia de la operación (paquete 7.1.4) con mentoría de seis meses (8.5.2). & Arquitecto de Solución & Comité de Arquitectura; Comité de Operación \\
  Planificador de Rutas & Alta / muy alto & Co-diseñar: talleres de reglas en los meses 1 a 3 (paquete 1.2.2), validación de las rutas propuestas en cada iteración de M4 y en la marcha blanca. & Líder de Implantación & Sesiones individuales; taller quincenal de rutas \\
  Preventistas (62) & Alta / alto & Consultar y capacitar: prueban la captura sin señal en su ruta y se certifican antes de la ola. & Líder de Implantación & Jefe de frente en ruta; reunión de inicio de jornada \\
  Tripulación propia y peonetas (84) & Alta / alto & Consultar y capacitar: prueban la entrega digital y la rendición antes de la ola. & Líder de Implantación & Jefe de frente en andén; reunión de inicio de jornada \\
  Conductores externos (aprox. 160) & Alta / alto & Consultar y capacitar: incorporación en el andén al asignarse (paquete 7.1.3), con autenticación del conductor real. & Líder de Implantación & Andén; guía rápida impresa y en el terminal \\
  Preparadores (120 en Talca y 60 en Concepción) & Alta / alto & Consultar y capacitar: aprendizaje en el puesto de dos horas como máximo y capacitación continua del turno de noche (paquete 7.1.2). & Líder de Implantación & Tutor por turno; reunión de cambio de turno \\
  Sindicato de Choferes & Alta / alto & Informar y acordar: define finalidades y controles de la telemetría; acuerdo previo antes de usar GPS para control de jornada (paquete 5.4.2). & Jefe de Proyecto & Mesa de trabajo con acta \\
  Empresas transportistas (10) & Media / alto & Informar y acordar: firman el acuerdo operacional antes de la ola de reparto (paquete 5.4.1). & Jefe de Proyecto & Reunión por empresa; acuerdo firmado \\
  Canal Tradicional (11.600 clientes) & Muy alta / alto & Informar: conserva su forma de compra y pago, sin exigirle internet ni dispositivo. & Líder de Implantación & Preventista; comprobante impreso \\
  \textit{Food service} (2.100 clientes) & Muy alta / alto & Informar: valida sus ventanas de entrega y la evidencia térmica del viaje. & Líder de Implantación & Ejecutivo comercial; correo \\
  Canal Moderno (500 puntos) & Muy alta / alto & Informar y acordar: certificación del intercambio electrónico por cadena (paquetes 3.6.5 y 3.6.6). & Arquitecto de Solución & Mesa técnica por cadena \\
  Proveedores (180) & Media / alto & Informar: aviso anticipado del rotulado GS1-128 y del portal de consulta. & Líder de Implantación & Comunicado formal; portal \\
  Autoridad Sanitaria & Media / alto & Informar: recibe la evidencia de trazabilidad y temperatura cuando la solicita. & Líder de Datos & Canal formal a través de la Jefa de Calidad \\
\end{tablalafrox}

El registro muestra dónde se concentra el esfuerzo de gestión. Los tres actores de la dirección deciden, por lo que su participación se canaliza por el Comité Ejecutivo y no por reuniones de trabajo. Las cinco jefaturas y el planificador co-diseñan: su conocimiento es el insumo de los requisitos y su validación precede a cada iteración que los usa. El mayor riesgo de resistencia está en el personal de terreno y en el sindicato, que se atienden en su puesto, en su horario y con acuerdos firmados antes de cada ola. Los clientes, proveedores y la autoridad sanitaria sólo reciben información y, cuando corresponde, firman un acuerdo, porque el proyecto no les exige cambiar su forma de trabajar.

\section*{Anexo 6.B \textemdash{} Adquisiciones asociadas al proyecto}
\addcontentsline{toc}{section}{Anexo 6.B \textemdash{} Adquisiciones asociadas al proyecto}

Este anexo reúne los bienes y servicios necesarios para implementar y operar la solución, e identifica para cada adquisición la necesidad que atiende, quién es responsable de su compra o contratación, cuándo se requiere y qué evidencia permitirá verificar su cumplimiento. Esta información facilita la coordinación de las adquisiciones con el cronograma y con las actividades del proyecto.

\setcounter{table}{0}\renewcommand{\thetable}{6.B.\arabic{table}}

El detalle se presenta en la Tabla~\ref{tab:tabla-adquisiciones}.
\begin{tablalafrox}{Adquisiciones -- Fuente: Formulario T-14, fases 5 y 6; Formulario T-11; Subdocumento 6, sección 6.1.3}{tab:tabla-adquisiciones}%
  {F{0.25} Y Y}%
  {\cab{Bien o servicio} & \cab{Necesidad asociada} & \cab{Responsable de compra o contrato; responsable LafroX, fecha, paquetes y evidencia}}
  Equipos de uso: terminales de preventa, reparto y bodega, impresoras, terminales de pago y balanzas & Operación de las aplicaciones & CLIENTE, según la especificación 5.1.1. Arquitecto (especificación); SRE (recepción). Antes de cada ola. Paquetes 6.5 y 4.2.1. Actas 5.1.3 \\
  Sistema de cadena de frío: gateways y sensores de cámara (B-01 y B-02) y kit de los 28 camiones (gateway, UPS, módulo de sondas, sondas y simuladores, sensor independiente, termógrafo e Iridium Edge) & Registro, protección local y aviso de temperatura & CLIENTE: sensores de cámara (B-01) y kits, según la especificación 5.1.4, con reserva de componentes para 31 kits. LafroX: gateways de cámara (B-02), como infraestructura on-premise. Arquitecto (especificación); SRE (instalación y recepción). Compra antes del banco de pruebas, con cotización local y certificación Subtel aplicable. Paquetes 2.4, 4.1 y 6.5. Actas por camión \\
  Sala técnica de Talca: UPS, generador, climatización de precisión, detección y extinción, racks y gabinetes de borde & Recinto técnico del H3 & LafroX, según la especificación 5.1.2. Arquitecto (especificación); SRE (instalación y recepción técnica). Orden de compra mes 2; instalación mes 3; recepción de sala mes 4 según T-15. Fecha de recepción técnica no especificada aquí. Paquetes 6.1, 6.3, 6.6. Actas 5.1.3 y 6.1.5 \\
  Servidores, almacenamiento, firewalls, switches y mini-PC de borde & Cómputo y red de los sitios & LafroX, según la especificación 5.1.2. SRE (instalación y recepción técnica). Orden de compra mes 2; instalación mes 3; recepción de sala mes 4 según T-15. Fecha de recepción técnica no especificada aquí. Paquetes 6.3, 6.6. Actas 5.1.3 \\
  Obra civil de separación, piso técnico e instalaciones del recinto & Habilitación del recinto & CLIENTE: obra civil de separación (RT-06.06). LafroX: especificación y coordinación; piso técnico e instalaciones con especialistas. SRE. Mes 3; recepción de sala mes 4. Paquetes 6.1.1 a 6.1.4; acta 6.1.5 \\
  Servicios de AWS & Plataforma de nube y ambientes & LafroX (5.2.1). SRE. Antes del mes 4. Paquetes 3.1, 3.2. Cuentas y servicios activos \\
  Gestión de dispositivos y detección en endpoints & Enrolamiento y seguridad de los terminales & LafroX (5.2.2). SRE. Antes del enrolamiento. Paquetes 6.5, 7.1. Suscripciones activas \\
  Licencias de software de terceros & Productos de la arquitectura & A nombre del CLIENTE (5.2.3). Arquitecto. Antes de usar cada producto. Paquetes 3.1 a 3.6. Registro de licencias \\
  Fibra óptica de Telsur (Grupo Gtd) en Talca y Concepción & Enlace principal & LafroX (5.3.1). SRE. Antes del H3. Paquete 6.6. Contrato y fecha de instalación \\
  Ocho planes LTE: Entel en Talca y Concepción; Entel y Claro en cada cross-docking & Respaldo de enlace & LafroX (5.3.2). SRE. Antes del H3. Paquete 6.6. Contratos \\
  Starlink de los cinco sitios & Enlace principal de los cross-docking y tercer camino de los CD & LafroX (5.3.3). SRE. Antes del H3. Paquete 6.6. Contratos y equipos \\
  SIM M2M de Entel y Claro y servicio satelital Iridium de los camiones con frío & Aviso y curva de temperatura en ruta & LafroX (5.3.4): 62 SIM M2M y servicio Iridium SBD. SRE. Antes del piloto de cadena de frío. Paquete 6.6. Contratos \\
  Acuerdos con los diez transportistas & Uso de terminales y suplentes enrolados & LafroX con el CLIENTE (5.4.1). Implantación. Antes del H6. Paquetes 4.2.1, ola de reparto. Diez acuerdos firmados \\
  Acta con el sindicato & Terminales, GPS y cámaras & CLIENTE y sindicato (5.4.2). Implantación. Antes del H6. Paquete 4.2.1. Acta firmada \\
  Custodia de fuentes & Continuidad ante insolvencia o incumplimiento & LafroX (5.4.4). Jefe de proyecto. Antes del H4. Paquete 9.2. Contrato y primer depósito \\
  Analistas del SOC 24×7 propio & Monitoreo de seguridad desde el mes 13 & LafroX (8.1.5; RT-11.17). Encargado de Seguridad. Contratación antes del mes 13; operación desde el mes 13. SOC propio en el Centro de Operaciones de Valparaíso; dotación y procedimientos en el Subdocumento 4, Anexo 4-B \\
\end{tablalafrox}


La tabla refleja que las adquisiciones tienen responsables y momentos de necesidad distintos: algunas corresponden al CLIENTE y otras a LafroX, y varias deben coordinarse con hitos o paquetes específicos. Por ello, el seguimiento de fechas y evidencias de recepción o contratación es clave para anticipar atrasos que puedan afectar la implementación y la continuidad de la operación.
